% TOP_PROBLEM: {"id": 20003127, "title": "The differentiable diameter sphere problem: a three-gate obstruction", "category_id": 7, "category": {"id": 7, "name": "topology", "display_name": "Topology", "description": "Properties preserved under continuous deformations.", "slug": "topology", "order_index": 7, "created_at": "2026-07-31T15:26:25.671Z"}, "status": "partially_solved", "problem_number": "AIM-TOPOLOGY-0215", "set_id": 14}
% FIRST_POSTED: 2026-10-01
% ENTRY_KIND: statement_only
% UnsolvedMath ID 20003127; source catalogue code AIM-TOPOLOGY-0215
% !TeX program = lualatex
% Definition and sources only; this file is not an LLM solution attempt.
\documentclass[11pt,a4paper]{article}
\usepackage[margin=25mm]{geometry}
\usepackage{fontspec}
\setmainfont{lmroman10-regular.otf}[BoldFont=lmroman10-bold.otf,
 ItalicFont=lmroman10-italic.otf,BoldItalicFont=lmroman10-bolditalic.otf]
\usepackage{amsmath,amssymb,mathtools}
\usepackage{unicode-math}
\setmathfont{latinmodern-math.otf}
\AtBeginDocument{\let\setminus\smallsetminus}
\usepackage{xurl,hyperref}
\hypersetup{colorlinks=true,urlcolor=blue}
\setcounter{secnumdepth}{0}
\setlength{\parindent}{0pt}
\setlength{\parskip}{5pt}
\setlength{\emergencystretch}{3em}
\begin{document}
\section*{The differentiable diameter sphere problem: a three-gate obstruction}\label{20003127}
{\small UnsolvedMath ID 20003127 · AIM-TOPOLOGY-0215 · Topology · AIM Workshop Problem Lists}
\subsection{Definitions and mathematical statement}
Let $(M,g)$ be a closed connected smooth Riemannian manifold of dimension $n\geq2$.
Closed means compact without boundary.
For a two-dimensional linear subspace $\sigma\subseteq T_xM$, let $\sec_g(\sigma)$ be its sectional curvature.
Write $d_g$ for geodesic distance and
$\operatorname{diam}(M,g)=\sup_{x,y\in M}d_g(x,y)$.
The curvature normalization $\sec_g\geq1$ means that every tangent two-plane at every point has curvature at least one.

The differentiable diameter question asks whether
\[
 \sec_g\geq1,\qquad
 \operatorname{diam}(M,g)>\frac\pi2
 \quad\Longrightarrow\quad
 M\text{ is diffeomorphic to }S^n.
\]
Here $S^n$ carries its standard smooth structure, and a diffeomorphism is a bijective smooth map with smooth inverse.
The conclusion concerns the smooth type of the manifold, not whether its given metric is round.
No upper sectional-curvature bound or pinching ratio is part of the hypothesis.
The strict diameter inequality is retained; the endpoint $\pi/2$ allows different familiar topological types and is a separate rigidity problem.
The classical diameter theorem provides the topological sphere conclusion under these hypotheses; the issue is whether an exotic smooth sphere could still occur.
The assertion is quantified over every dimension and every metric satisfying the displayed normalization, with no symmetry or simple-connectivity assumption added.
Equivalently, it asks whether any nonstandard smooth sphere is excluded from this entire normalized diameter range.
\subsection{Short English statement}
Does curvature at least one and diameter greater than $\pi/2$ force the standard smooth sphere, rather than only its topological type?
\subsection{Sources}
\begin{itemize}
\item[S1] ULAM.AI and the UnsolvedMath Contributors, supplied problems.json, record 20003127 (AIM-TOPOLOGY-0215), AIM Workshop Problem Lists. Catalogue identity and source transcription; CC BY 4.0.
\url{https://huggingface.co/datasets/ulamai/UnsolvedMath}.
\item[S2] American Institute of Mathematics, Manifolds with nonnegative sectional curvature, item 2. Original workshop question underlying AIM-TOPOLOGY-0215.
\url{https://aimath.org/WWN/nnsectcurvature/nnsectcurvature.pdf}.
\end{itemize}
\end{document}
